\documentclass[10pt,a4paper]{amsart}
\usepackage[margin=19mm]{geometry}
\usepackage{amsmath,amssymb,mathtools}
\usepackage[scr=rsfs,cal=euler]{mathalfa}
\usepackage{xfrac}
\usepackage[unicode,hidelinks,bookmarks=false]{hyperref}
\newcommand{\ie}{i.e.\ }
\newcommand{\cf}{cf.\ }
\newcommand{\resp}{respectively\ }
\newcommand{\Subsection}{\subsection}
\providecommand{\nobreakdash}{\nobreak\hbox{-}}
\setlength{\emergencystretch}{2em}
\multlinegap0pt
\usepackage{bm}
\usepackage{bbm}
\usepackage{dsfont}
\usepackage{xspace}
\usepackage{cancel}
\usepackage{mathtools}
\usepackage{amsthm}
\makeatletter
\@ifundefined{theorem}{\newtheorem{theorem}{Theorem}}{}
\makeatother
\title[A recursive formula for the $n^\text{th}$ survival function ]{A recursive formula for the $n^\text{th}$ survival function and the $n^\text{th}$ first passage time distribution for jump and diffusion processes. Applications to the pricing of $n^\text{th}$-to-default CDS}
\author{Alessio Lapolla}
\date{Source: arXiv:2509.02347v2 (CC BY 4.0)}
\begin{document}
\maketitle

\noindent\emph{Source and licence.} A recursive formula for the $n^\text{th}$ survival function and the $n^\text{th}$ first passage time distribution for jump and diffusion processes. Applications to the pricing of $n^\text{th}$-to-default CDS. Version arXiv:2509.02347v2 (CC BY 4.0), distributed under the Creative Commons Attribution 4.0 International licence, as recorded in the arXiv licence metadata for this exact version and shown on the abstract page https://arxiv.org/abs/2509.02347v2. Full rights notice: licenses/arxiv-2509.02347-CC-BY-4.0.txt.\par\smallskip

\noindent\textit{Editorial scope.} Counted unit: the Theorem whose source label is \texttt{theo:two} in the article (source0.tex lines 81--92, source label \texttt{theo:two}), delivered together with the article's own complete original proof of it (source0.tex lines 93--119, one \texttt{proof} environment of 27 lines). No mathematics is written, repaired or re-derived: statement and proof are the article's own text line for line. Required context retained uncounted: none. Every definition, display and label of those lines is retained unchanged. Omitted from this package and not redistributed: 1--80, 120--501. Nothing inside a retained excerpt is omitted or altered: every retained body is delivered exactly as the article's lines are written, so no omission receipt of any kind accompanies this package. Citations: the retained excerpts carry no citation command at all, so no bibliography entry of the article is transcribed and no reference is redistributed. Reference adaptation: none was needed and none was made; every reference inside the delivered unit resolves inside it, so no unresolved reference or citation is produced. Printed slips kept literal: no printed word, symbol, hypothesis or number of the counted statement or of its proof was altered. Layout and class: the article's own class is \texttt{amsart} with packages including inputenc, graphicx, babel, setspace, amsmath, amssymb, amsfonts, booktabs, amsthm, gnuplot-lua-tikz, mathtools, hyperref, tikz; the wrapper is the lane's pinned \texttt{amsart} editorial wrapper at 10pt on A4, which restates the theorem-like environments the retained text needs and adds the article's own macro definitions that the retained text needs (none), with extra wrapper packages (bm, bbm, dsfont, xspace, cancel, mathtools). The delivered numbering is the unit's own. Assets: no retained span contains a figure environment, a graphics-inclusion command, a PDF-page inclusion, a table or a TikZ picture; no image file is copied into this package, the package declares no asset dependency and no image file is redistributed with it. Licence: version arXiv:2509.02347v2 of this article is distributed under the Creative Commons Attribution 4.0 International licence, as recorded in the arXiv licence metadata for this exact version and shown on the abstract page https://arxiv.org/abs/2509.02347v2. A full rights notice is delivered at licenses/arxiv-2509.02347-CC-BY-4.0.txt. Characters: every retained excerpt is pure ASCII, so the delivered engine, XeLaTeX with its default Latin Modern text font, reports no missing character.
  \begin{theorem}
    \label{theo:two}
    Assuming that all $P$s are probability density function and that $S_1(t)$ and $S_2(t)$ are absolutely continuous, then the last particle survival function $S^1(t)$ can be computed using the following formula:
    \begin{equation}
      \label{eq:S1_general}
      \begin{aligned}
        &S^1(t)=S^2(t)+\int_Rdx\int_0^tds \int_Rdy P^1_1(x,t|y,s)P^2_1(y,s|X_2\in \partial R,s)F_{2}(s)+\\
        &\int_Rdx\int_0^tds \int_Rdy P^1_2(x,t|y,s)P^2_2(y,s|X_1\in \partial R,s)F_{1}(s),
      \end{aligned}
    \end{equation}
    and $F^1(t)=-\frac{dS^1(t)}{dt}$.
  \end{theorem}

  \begin{proof}
    By definition the law of either coordinate $1$ or $2$ is described by $P^2_i(x,t|x_0,s)$ for $s<t\leq\tau_m$ and $i=1,2$; while the law of the last coordinate remaining by $P^1_1(x,t|x_0,s)$ if $\tau_2<\tau_1$ or $P^1_2(x,t|x_0,s)$ if $\tau_2>\tau_1$ for $t>s>\tau_m$. Therefore the probability of survival of the last particle is composed by two mutually exclusive and complementary events:
    \begin{equation}
      S^1(t)=\mathbb{P}(t<\tau_M)=\mathbb{P}(t\leq\tau_m)+\mathbb{P}(\tau_m<t<\tau_M)=S^2(t)+\mathbb{P}(\tau_m<t<\tau_M),
    \end{equation}
    where the first addend encompasses also the case $\mathbb{P}(\tau_1=\tau_2)$, even when this probability is larger than zero. The second addend is in turn formed by two mutually exclusive events, and thanks to the law of total probability:
    \begin{equation}
      \label{eq:total_prob}
      \begin{aligned}
        \mathbb{P}(\tau_m<t<\tau_M)=&\mathbb{P}(\tau_m<t<\tau_M|\tau_2=\tau_m)\mathbb{P}(\tau_2=\tau_m)+\\
        &\mathbb{P}(\tau_m<t<\tau_M|\tau_1=\tau_m)\mathbb{P}(\tau_1=\tau_m).
      \end{aligned}
    \end{equation}
    Now, since the first killing happened at a certain time $s<\tau_m<t$ we have that
    \begin{equation}
      \mathbb{P}(\tau_2=\tau_m)=\mathbb{P}(t>\tau_2)=1-\mathbb{P}(t\leq\tau_2)=1-S_2(t)
    \end{equation}
    is the default probability of $X_2$, and
    \begin{equation}
      \begin{aligned}
        &\mathbb{P}(\tau_m<t<\tau_M|\tau_2=\tau_m)=\frac{\int_0^tds \mathbb{P}(s<t<\tau_M|\tau_2=s)F_{2}(s)}{1-\mathbb{P}(t\leq\tau_2)}=\\
        &\frac{\int_0^tds \int_Rdy \mathbb{P}(t<\tau_M|y,s)P^2_1(y,s|X_2 \in \partial R,s)F_{2}(s)}{1-\mathbb{P}(t\leq\tau_2)}=\\
        &\frac{\int_Rdx\int_0^tds \int_Rdy P^1_1(x,t|y,s)P^2_1(y,s|X_1 \in \partial R,s)F_{2}(s)}{1-\mathbb{P}(t\leq\tau_2)};
      \end{aligned}
    \end{equation}
    where we computed the conditional average over the killing time of $X_2$ in the first equality, used the Markov property in the second, and used the definition of survival function in the third. The same applies\emph{ mutatis mutandis} to the second added of Eq.~\eqref{eq:total_prob}. Thus, putting all together, we finally obtain Eq.~\eqref{eq:S1_general}.
  \end{proof}

\end{document}
